\chapter{Segitiga Siku-siku dan Kosinus}\label{ch:g8:cosine}

Dua kenyataan yang indah mengikat \omterm{def:g6:shapes:triangles}{segitiga siku-siku} pada lingkaran:
\omterm{def:g6:shapes:triangles}{segitiga siku-siku} muat persis dalam setengah lingkaran, dengan sisi
miringnya menjadi \omterm{def:g4:geometry:circle}{diameter}. Dan satu bilangan, yaitu \emph{kosinus}
sebuah sudut, menyandikan bentuk setiap \omterm{def:g6:shapes:triangles}{segitiga siku-siku} yang
bersudut itu --- perbandingan trigonometri yang pertama, mendahului
saudaranya sinus dan tangen (\cref{ch:g9:trig}).

\section{Segitiga siku-siku dan lingkarannya}

\begin{theorem}[Teorema lingkaran]\label{thm:g8:cosine:circle}
\begin{enumerate}
  \item Kalau segitiga $ABC$ siku-siku di $A$, maka $A$ terletak pada
        lingkaran yang \omterm{def:g4:geometry:circle}{diameternya} sisi miringnya $[BC]$.
  \item Sebaliknya, kalau $A$ terletak pada lingkaran berdiameter
        $[BC]$ (dengan $A \neq B, C$), maka segitiga $ABC$ siku-siku
        di $A$.
\end{enumerate}
Setara dengan itu: pada \omterm{def:g6:shapes:triangles}{segitiga siku-siku}, titik tengah sisi miringnya
berjarak sama dari ketiga \omterm{def:g5:solids:def}{titik sudutnya} --- \omterm{pb:g8:midpoints:1}{garis berat} dari \omterm{def:g3:shapes:rightangle}{sudut
siku-sikunya} berukuran setengah sisi miringnya.
\end{theorem}

\begin{proof}[Bukti untuk 1]
Misalkan $O$ titik tengah $[BC]$ dan $D$ titik \omterm{def:g7:central:def}{simetris} $A$ terhadap
$O$. Diagonal $ABDC$ berpotongan di titik tengah bersamanya $O$, jadi
$ABDC$ adalah \omterm{def:g7:central:parallelogram}{jajargenjang} (\cref{def:g7:central:parallelogram}) ---
dengan \omterm{def:g3:shapes:rightangle}{sudut siku-siku} di $A$: jadi ia persegi panjang. Diagonal
persegi panjang sama panjang, jadi $OA = \frac{AD}{2} = \frac{BC}{2}$:
titik $A$ berjarak $\frac{BC}{2}$ dari $O$, yaitu terletak pada
lingkaran berdiameter $[BC]$. (Bagian 2 dibuktikan dengan menjalankan
alasannya mundur.)
\end{proof}

\begin{omfigure}
\begin{tikzpicture}[scale=1.15]
  \coordinate (B) at (-1.8,0);
  \coordinate (C) at (1.8,0);
  \coordinate (O) at (0,0);
  \coordinate (A) at (130:1.8);
  \coordinate (A2) at (55:1.8);
  \draw[thick] (O) circle (1.8);
  \draw[thick] (B) -- (C);
  \draw[thick, omDef] (B) -- (A) -- (C);
  \draw[thick, omExo] (B) -- (A2) -- (C);
  \draw[omThm, thick] (O) -- (A);
  \fill (B) circle (1.5pt) node[below left, font=\small] {$B$};
  \fill (C) circle (1.5pt) node[below right, font=\small] {$C$};
  \fill (O) circle (1.5pt) node[below, font=\small] {$O$};
  \fill (A) circle (1.5pt) node[above left, font=\small] {$A$};
  \fill (A2) circle (1.5pt) node[above right, font=\small] {$A'$};
  \draw[thin] (A) ++(0.14,-0.20) -- ++(0.14,0.10) -- ++(-0.14,0.20);
\end{tikzpicture}

{\small Di mana pun $A$ duduk pada lingkarannya, sudut $\widehat{BAC}$
siku-siku --- yaitu dugaan pada \cref{exo:g6:shapes:10}, yang kini
menjadi teorema. \omterm{pb:g8:midpoints:1}{Garis berat} merahnya $[OA]$ adalah \omterm{def:g3:shapes:shapes}{jari-jari}: yaitu
setengah sisi miringnya.}
\end{omfigure}

\begin{example}\label{ex:g8:cosine:circle}
Sebuah segitiga punya sisi miring $10$ cm. Tanpa mengetahui apa pun
lagi, \omterm{pb:g8:midpoints:1}{garis berat} dari \omterm{def:g3:shapes:rightangle}{sudut siku-sikunya} berukuran $5$ cm, dan
\omterm{pb:g6:symmetry:1}{lingkaran luar} segitiganya berjari-jari $5$ cm, berpusat di
titik tengah sisi miringnya.
\end{example}

\section{Kosinus sebuah sudut lancip}

\begin{definition}[Kosinus]\label{def:g8:cosine:def}
Pada \omterm{def:g6:shapes:triangles}{segitiga siku-siku}, untuk sudut lancip $\theta$:
\[
\cos\theta = \frac{\text{sisi samping } \theta}{\text{sisi miring}}
\]
--- yaitu sisi siku-siku yang menyentuh $\theta$, dibagi sisi
miringnya. Perbandingan ini hanya bergantung pada sudutnya, bukan pada
besar segitiganya: semua \omterm{def:g6:shapes:triangles}{segitiga siku-siku} dengan sudut lancip yang
sama adalah perbesaran satu sama lain, dan perbesaran mempertahankan
perbandingan panjang (\cref{ch:g8:midpoints} memulai kisah ini;
\cref{ch:g9:thales} menyelesaikannya).
\end{definition}

\begin{omfigure}
\begin{tikzpicture}[scale=1.05]
  \coordinate (B) at (0,0);
  \coordinate (A) at (4.6,0);
  \coordinate (C) at (4.6,2.3);
  \coordinate (A1) at (2.76,0);
  \coordinate (C1) at (2.76,1.38);
  \draw[thick] (B) -- (A) -- (C) -- cycle;
  \draw[omThm, thick] (A1) -- (C1);
  \draw[thin] (A) ++(-0.25,0) -- ++(0,0.25) -- ++(0.25,0);
  \draw[thin] (A1) ++(-0.22,0) -- ++(0,0.22) -- ++(0.22,0);
  \draw[omDef, ->] (0.85,0) arc (0:26.6:0.85);
  \node[omDef, font=\small] at (1.25,0.24) {$\theta$};
  \node[below left, font=\small] at (B) {$B$};
\end{tikzpicture}

{\small Dua \omterm{def:g6:shapes:triangles}{segitiga siku-siku} yang berbagi sudut $\theta$: yang kecil
adalah pengecilan yang besar, jadi
$\frac{\text{samping}}{\text{miring}}$ sama untuk keduanya ---
nilai bersama itulah $\cos\theta$.}
\end{omfigure}

\begin{proposition}[Nilai dan batas pertamanya]\label{prop:g8:cosine:bounds}
Untuk setiap sudut lancip $\theta$: $0 < \cos\theta < 1$, dan
kosinusnya \emph{turun} ketika sudutnya membuka: sudut yang lebih
lebar punya kosinus yang lebih kecil. Patokannya:
$\cos 0^\circ = 1$, $\cos 60^\circ = \frac12$, $\cos 90^\circ = 0$
(nilai ekstremnya bersesuaian dengan segitiga yang memipih).
\end{proposition}
\admitted

\begin{method}[Memakai kosinus]\label{met:g8:cosine:use}
Pada \omterm{def:g6:shapes:triangles}{segitiga siku-siku}, ketika besaran yang diketahui dan yang dicari
adalah sudut lancip, sisi sampingnya, dan sisi miringnya:
\begin{enumerate}
  \item tulislah definisinya:
        $\cos\theta = \frac{\text{samping}}{\text{miring}}$ dengan
        nilai yang diketahui pada tempatnya;
  \item selesaikan \omterm{def:g8:equations:def}{persamaan} kecilnya untuk yang belum diketahui
        (dengan mengalikan atau membagi);
  \item untuk \emph{sudut} yang belum diketahui, terapkan
        $\cos^{-1}$ pada kalkulator terhadap perbandingan yang
        dihitung;
  \item pemeriksaan akal sehat: panjangnya harus keluar lebih pendek
        daripada sisi miringnya; sudutnya benar-benar antara
        $0^\circ$ dan $90^\circ$.
\end{enumerate}
\end{method}

\begin{example}[Mencari sebuah sisi]\label{ex:g8:cosine:side}
Jalur landai sepanjang $6$ m membentuk sudut $20^\circ$ dengan tanah
yang mendatar. Jarak mendatar yang ditempuhnya (sisi samping
$20^\circ$, dengan sisi miring $6$ m):
\[
\cos 20^\circ = \frac{d}{6}
\quad\Longrightarrow\quad
d = 6 \cos 20^\circ \approx 6 \times 0.940 \approx 5.6 \text{ m}.
\]
\end{example}

\begin{example}[Mencari sebuah sudut]\label{ex:g8:cosine:angle}
Pada \omterm{def:g6:shapes:triangles}{segitiga siku-siku}, sisi samping sudut $\theta$ berukuran
$3.5$ dan sisi miringnya $5$:
\[
\cos\theta = \frac{3.5}{5} = 0.7,
\qquad
\theta = \cos^{-1}(0.7) \approx 45.6^\circ .
\]
\end{example}

\section{Latihan}

\begin{exercise}[$\star$]\label{exo:g8:cosine:1}
Sebuah \omterm{def:g6:shapes:triangles}{segitiga siku-siku} punya sisi miring $13$ cm. Berapa panjang
\omterm{pb:g8:midpoints:1}{garis berat} dari titik \omterm{def:g3:shapes:rightangle}{sudut siku-sikunya}? Berapa \omterm{def:g3:shapes:shapes}{jari-jari} lingkaran
yang melalui ketiga \omterm{def:g5:solids:def}{titik sudutnya}?
\end{exercise}

\begin{exercise}[$\star$]\label{exo:g8:cosine:2}
Gambarlah lingkaran berdiameter $8$ cm dengan \omterm{def:g4:geometry:circle}{diameter} $[BC]$,
pilihlah titik $A$ mana pun pada lingkarannya lalu gambarlah $ABC$.
Sudut yang mana yang siku-siku? Sebutkan teoremanya. Ukurlah $AB$ dan
$AC$ lalu periksalah Pythagoras.
\end{exercise}

\begin{exercise}[$\star$]\label{exo:g8:cosine:3}
Pada segitiga $DEF$ yang siku-siku di $D$, sebutkan sisi samping
sudut $\widehat E$, sisi depannya, dan sisi miringnya. Tulislah
$\cos \widehat E$ sebagai sebuah perbandingan.
\end{exercise}

\begin{exercise}[$\star$]\label{exo:g8:cosine:4}
Hitunglah besaran yang hilang ($\cos 35^\circ \approx 0.819$,
$\cos 50^\circ \approx 0.643$):
\begin{enumerate}
  \item sisi miring $10$, sudut $35^\circ$: berapa sisi sampingnya?
  \item sisi samping $6$, sudut $50^\circ$: berapa sisi miringnya?
\end{enumerate}
\end{exercise}

\begin{exercise}[$\star$]\label{exo:g8:cosine:5}
Pada sebuah \omterm{def:g6:shapes:triangles}{segitiga siku-siku}, sisi samping $\theta$ berukuran $8$
dan sisi miringnya $10$. Hitunglah $\cos\theta$, lalu $\theta$
($\cos^{-1}(0.8) \approx 36.9^\circ$).
\end{exercise}

\begin{exercise}[$\star$]\label{exo:g8:cosine:6}
Jelaskan mengapa $\cos\theta$ tidak pernah dapat bernilai $1.2$ untuk
sudut lancip sebuah \omterm{def:g6:shapes:triangles}{segitiga siku-siku}.
\end{exercise}

\begin{exercise}[$\star\star$]\label{exo:g8:cosine:7}
Sebuah luncuran tali sepanjang $25$ m menurun dari sebuah anjungan ke
tanah, dengan membentuk sudut $12^\circ$ terhadap arah mendatar
($\cos 12^\circ \approx 0.978$). Berapa jarak mendatar yang direntangnya?
\end{exercise}

\begin{exercise}[$\star\star$]\label{exo:g8:cosine:8}
Urutkan tanpa kalkulator: $\cos 20^\circ$, $\cos 70^\circ$,
$\cos 45^\circ$ (pakailah \cref{prop:g8:cosine:bounds}). Lalu
periksalah dengan kalkulator.
\end{exercise}

\begin{exercise}[$\star\star$]\label{exo:g8:cosine:9}
Tangga sepanjang $L$ bersandar pada tembok dengan sudut $\theta$
antara tangga dan \emph{tanah}. Kakinya $1.2$ m dari tembok
dan $\theta = 68^\circ$ ($\cos 68^\circ \approx 0.375$). Hitunglah
$L$, lalu tinggi yang tercapainya (dengan Pythagoras).
\end{exercise}

\begin{exercise}[$\star\star$]\label{exo:g8:cosine:10}
$[BC]$ adalah \omterm{def:g4:geometry:circle}{diameter} sebuah lingkaran berpusat $O$ dan berjari-jari
$4.5$ cm, dan $A$ adalah titik pada lingkarannya dengan $AB = 5.4$ cm.
\begin{enumerate}
  \item Mengapa $ABC$ siku-siku di $A$?
  \item Hitunglah $AC$, lalu $\cos \widehat B$.
\end{enumerate}
\end{exercise}

\begin{exercise}[$\star\star\star$]\label{exo:g8:cosine:11}
Dengan memakai setengah \omterm{def:g6:shapes:triangles}{segitiga sama sisi} bersisi $1$ (yang dipotong
sepanjang sebuah tinggi), berilah alasan untuk patokan
$\cos 60^\circ = \frac12$ pada \cref{prop:g8:cosine:bounds}. (Di mana
kaki tingginya jatuh pada alasnya? Versi trigonometri lengkapnya
adalah \cref{ex:g9:trig:special}.)
\end{exercise}

\section{Soal: Hubungan Euclid, dan Pythagoras sekali lagi}

\begin{problem}[{Soal akhir pekan --- garis tinggi segitiga
siku-siku: $AB^2 = BH \times BC$, $AH^2 = BH \times HC$, bukti kedua
Pythagoras, dan mesin untuk melukis akar kuadrat}]
\label{pb:g8:cosine:1}
Turunkanlah, dari \omterm{def:g3:shapes:rightangle}{sudut siku-siku} sebuah \omterm{def:g6:shapes:triangles}{segitiga siku-siku}, \omterm{def:g6:lines:perp}{garis
tegak lurus} ke sisi miringnya: \omterm{def:g6:lines:objects}{ruas garis} pendek ini --- yaitu
\emph{garis tingginya} --- memenuhi hubungan yang begitu berguna
sehingga Euclid menaruhnya di jantung karyanya \emph{Elements}. Pada
soal ini sifat yang mendefinisikan kosinus, ``perbandingannya hanya
bergantung pada sudutnya''
(\cref{def:g8:cosine:def}), membuktikan semuanya; di sepanjang jalannya
kamu akan membuktikan ulang teorema Pythagoras lewat jalan yang sama
sekali berbeda, lalu berakhir dengan mesin penggaris dan jangka yang
melukis $\sqrt2$, $\sqrt5$, $\sqrt{n}$ untuk setiap bilangan bulat $n$.

Sepanjang soal ini, $ABC$ adalah segitiga yang siku-siku di $A$, dan
$H$ adalah kaki \omterm{def:g6:lines:perp}{garis tegak lurus} dari $A$ ke sisi miringnya $[BC]$,
sehingga $H$ terletak antara $B$ dan $C$ dan $AH \perp BC$.

\medskip
\noindent\textbf{Bagian I --- Satu sudut, dua segitiga.}
\begin{enumerate}
  \item Dengan memakai \omterm{def:g1:addition:def}{jumlah} sudut segitiga
        (\cref{thm:g7:triangles:sum}) pada $ABH$ dan pada $ABC$,
        tunjukkan bahwa $\widehat{BAH} = \widehat{ACB}$: garis
        tingginya memotong \omterm{def:g6:shapes:triangles}{segitiga siku-sikunya} menjadi dua segitiga
        yang lebih kecil yang membawa \emph{sudut yang sama} seperti
        aslinya.
  \item Segitiga $ABC$ (yang siku-siku di $A$) dan $ABH$
        (yang siku-siku di $H$) berbagi sudut $\widehat B$. Tulislah
        $\cos\widehat B$ pada masing-masingnya, lalu simpulkan dari
        \cref{def:g8:cosine:def} bahwa
        \[
        AB^2 = BH \times BC .
        \]
        (Untuk berpindah dari perbandingan yang sama ke \omterm{def:g2:mult:def}{hasil kali}
        yang sama, kalikan kedua ruasnya dengan kedua \omterm{def:g4:fractions:def}{penyebutnya},
        \cref{thm:g8:equations:balance}.)
  \item Nyatakan lalu buktikan hubungan kembarnya di \omterm{def:g5:solids:def}{titik sudut} $C$:
        \[
        AC^2 = CH \times CB .
        \]
  \item Ambillah segitiga $3$--$4$--$5$: $AB = 3$, $AC = 4$,
        $BC = 5$. Hitunglah $BH$ dan $CH$ dari pertanyaan 2 dan 3,
        lalu periksalah bahwa $BH + HC = BC$.
  \item Secara umum, nyatakan $BH$ dan $CH$ sebagai \omterm{def:g6:fractions:def}{pecahan} yang
        hanya memuat ketiga sisinya. Dalam arti apa kaki $H$
        membelah sisi miringnya ``\omterm{def:g6:propdata:def}{sebanding} dengan kuadrat kedua
        sisi siku-sikunya''?
\end{enumerate}

\medskip
\noindent\textbf{Bagian II --- Pythagoras lagi, dan garis tingginya.}
\begin{enumerate}[resume]
  \item Jumlahkan hubungan pada pertanyaan 2 dan 3 lalu pakailah
        $BH + HC = BC$ untuk memperoleh
        \[
        AB^2 + AC^2 = BC^2 :
        \]
        yaitu teorema Pythagoras (\cref{thm:g8:pythagoras:direct}),
        yang dibuktikan ulang --- tanpa teka-teki luas sama sekali
        (bandingkan \cref{exo:g8:pythagoras:11}).
  \item Terapkan Pythagoras pada segitiga kecil $ABH$ untuk menulis
        $AH^2 = AB^2 - BH^2$, gantilah $AB^2$ dengan
        $BH \times BC$, lalu faktorkan $BH$ untuk membuktikan
        \emph{hubungan garis tinggi Euclid}:
        \[
        AH^2 = BH \times HC .
        \]
  \item Hitunglah luas $ABC$ dengan dua cara --- sisi siku-sikunya
        sebagai alas dan tinggi, lalu sisi miringnya sebagai alas
        (\cref{thm:g7:areas:triangle}) --- lalu simpulkan hubungan
        yang ketiga:
        \[
        AB \times AC = BC \times AH .
        \]
  \item Kembali ke segitiga $3$--$4$--$5$: hitunglah $AH$ dengan
        pertanyaan 8, lalu periksalah hubungan garis tinggi pada
        pertanyaan 7 secara angka, memakai nilai $BH$ dan $CH$ yang
        ditemukan pada pertanyaan 4.
  \item Kaki garis tingginya membelah sisi miring suatu \omterm{def:g6:shapes:triangles}{segitiga
        siku-siku} menjadi ruas $BH = 2$ cm dan $HC = 8$ cm.
        Hitunglah garis tingginya $AH$.
\end{enumerate}

\medskip
\noindent\textbf{Bagian III --- Mesin untuk akar kuadrat.}
Gambarlah \omterm{def:g6:lines:objects}{ruas garis} $[BC]$ yang tersusun dari dua potongan yang
ditata ujung ke ujung: $BH = p$ dan $HC = q$. Gambarlah setengah
lingkaran berdiameter $[BC]$, lalu misalkan $A$ titik tempat \omterm{def:g6:lines:perp}{garis
tegak lurus} terhadap $(BC)$ di $H$ memotongnya.
\begin{enumerate}[resume]
  \item Dengan memakai \cref{thm:g8:cosine:circle}, jelaskan mengapa
        segitiga $ABC$ siku-siku di $A$ --- sehingga pertanyaan 7
        berlaku dan
        \[
        AH^2 = p \times q .
        \]
        (Panjang $AH$ disebut \emph{rata-rata
        geometri}\index{rata-rata geometri} dari $p$ dan $q$.)
  \item Ambillah $p = 1$ dan $q = 2$. Berapa $AH^2$? Panjang terkenal
        yang mana pada \cref{ex:g8:pythagoras:diagonal} yang baru
        saja dilukis gambar penggaris dan jangkamu?
  \item Uraikan resep yang melukis \omterm{def:g6:lines:objects}{ruas garis} sepanjang
        $\sqrt n$ untuk bilangan bulat $n \geq 1$ mana pun, lalu
        sebutkan pilihan $p$ dan $q$ yang melukis $\sqrt 5$.
  \item Misalkan $O$ pusat setengah lingkarannya. Jelaskan mengapa
        $AH$ tidak pernah dapat melampaui \omterm{def:g3:shapes:shapes}{jari-jarinya}
        $\frac{p + q}{2}$, dan untuk letak $H$ yang mana keduanya
        sama.
  \item Simpulkan ketaksamaannya: untuk semua bilangan positif $p$
        dan $q$,
        \[
        p \times q \leq \left(\frac{p + q}{2}\right)^2,
        \]
        dengan kesamaan persis ketika $p = q$. Lalu buktikan ulang
        tanpa geometri sama sekali: jabarkan
        $\left(\frac{p+q}{2}\right)^2 - p q$ lalu kenalilah sebuah
        kesamaan istimewa (\cref{pb:g8:equations:1}).

\end{enumerate}
\end{problem}
